% latex-oef3.tex
%
\documentclass[11pt,a4paper]{amsart}
\usepackage{amsmath,amssymb,amsthm}
\usepackage[utf8]{inputenc}
\usepackage[dutch]{babel}
\usepackage{enumitem}
\usepackage{graphicx}
\begin{document}
Goniometrie
Relaties tussen hoeken
Met behulp van de goniometrische cirkel worden de volgende relaties zichtbaar:
De tangens van een hoek is gedefinieerd als de verhouding tussen
 dus:
zodat
Uit de stelling van Pythagoras volgt:
Secans (sec), cosecans (csc) en cotangens (cot) zijn de reciproque functies van
respectievelijk: de cosinus, sinus en tangens.
Met de basisrelaties kan steeds een van de functies in een andere worden
uitgedrukt, zij het slechts voor een rechthoekige driehoek. De kunst bij
goniometrie is dan ook vaak om een willekeurige driehoek of veelhoek op te 
delen in rechthoekige driehoeken, zodat de basisrelaties toegepast kunnen 
worden. De volgende 6 formules gelden voor  tussen 0 en  radialen.

Verdere omrekenregels
De som- en verschilregels:
en zo kunnen we nog uren doorgaan!
 
\end{document}
